\subsection{Arithmetic and geometric sequences}

Two elementary families already show different kinds of repeated change.

\begin{definitionbox}[Arithmetic sequence]
An arithmetic sequence has a constant difference:
\[
a_{n+1}-a_n=d.
\]
If the first term is \(a_1\), then
\[
a_n=a_1+(n-1)d.
\]
\end{definitionbox}

\begin{examplebox}[Constant additive change]
For \(a_1=4\) and \(d=3\),
\[
4,7,10,13,\ldots
\]
and
\[
a_n=4+3(n-1)=3n+1.
\]
The index increases by one while the value increases by three.
\end{examplebox}

\begin{definitionbox}[Geometric sequence]
A geometric sequence has a constant ratio:
\[
\frac{a_{n+1}}{a_n}=q
\]
when the quotient is defined. Its explicit form is
\[
a_n=a_1q^{n-1}.
\]
\end{definitionbox}

\begin{examplebox}[Constant multiplicative change]
For \(a_1=8\) and \(q=\tfrac12\),
\[
8,4,2,1,\frac12,\ldots
\]
and
\[
a_n=8\left(\frac12\right)^{n-1}.
\]
This sequence models repeated proportional decay.
\end{examplebox}

\begin{interpretationbox}[Two models of change]
Arithmetic sequences describe repeated addition. Geometric sequences describe repeated multiplication. This distinction later becomes the contrast between linear and exponential behavior.
\end{interpretationbox}
